\documentclass[compacto,10pt]{pubciencias-notas}

% -- Paquetes adicionales
\usepackage{comandosEPN}

% -- Datos del libro
\autor{Departamento de Formación Básica}
\materia{Álgebra Lineal}
\nota{Semana 06 Clase 01}
\fecha{Semestre 2025-A}

\logouno[1.1cm]{Logo_EPN.eps}
\logodos[1.7cm]{Logo_DFB.eps}

\definecolor{colordef}{cmyk}{0.81,0.62,0.00,0.22}

\begin{document}

\encabezado

\vspace*{-8mm}

\setcounter{section}{2}
\section{Determinantes}
\setcounter{subsection}{1}
\subsection{Resolución de sistemas de ecuaciones}

\begin{defi}[Matriz de cofactores]
Sea $A\in\Mat[\K]{n}{n}$. La matriz de cofactores de $A$, que se denota por $\cof(A)$, es la matriz de orden $n \times n$ que está formada por los cofactores de $A$, es decir:
\[
\cof(A) = (C_{ij}) = 
\begin{pmatrix}
C_{11}&C_{12}&\ldots&C_{1n}\\
C_{21}&C_{22}&\ldots&C_{2n}\\
\vdots&\vdots&\ddots&\vdots\\
C_{n1}&C_{n2}&\ldots&C_{nn}
\end{pmatrix}.
\]
donde, $C_{ij} = (-1)^{i+j} \det (A_{ij})$.
\end{defi}

\begin{ejem}\label{matrizcofactores}
    Sea $A = 
    \begin{pmatrix}
    0 & 1 & 3 \\
    2 & -1 & 0 \\
    -1 & 2 & -1
    \end{pmatrix}$, determinemos sus cofactores correspondientes a la primera columna.
    \begin{itemize}
    \item
        $C_{11} = (-1)^{2} 
        \begin{vmatrix}
            -1 & 0 \\
            2 & -1
        \end{vmatrix} = 1$.
        
    \item
        $C_{21} = (-1)^{3}
        \begin{vmatrix}
            1 & 3 \\
            2 & -1
        \end{vmatrix} = 7$.
        
    \item
        $C_{31} = (-1)^{4}
        \begin{vmatrix}
            1 & 3 \\
            -1 & 0
        \end{vmatrix} = 3$.
    \end{itemize}
\end{ejem}

\begin{ejer}
    Verifique que la matriz de los cofactores de la matriz $A$ del ejemplo \ref{matrizcofactores} es: \[ \cof(A) = 
    \begin{pmatrix}
            1 & 2 & 3\\
            7 & 3 & -1 \\
            3 & 6 & -2
        \end{pmatrix}
    \]
\end{ejer}

\begin{defi}[Matriz adjunta]
Sea $A\in\Mat[\K]{n}{n}$. La matriz adjunta de $A$, que se denota por $\adj(A)$, es la matriz de orden $n \times n$ igual a la transpuesta de la matriz de cofactores de $A$, es decir:
\[
    \adj(A) = \cof(A)^{T}.
\]
\end{defi}

\begin{ejem}
    Consideramos la matriz $A$, planteada en el ejemplo \ref{matrizcofactores}. Se tiene:
    \[
        \adj(A) = \cof(A)^{T} =
        \begin{pmatrix}
            1 & 7 & 3 \\
            2 & 3 & 6 \\
            3 & -1 & -2
        \end{pmatrix}
    \]
\end{ejem}

\begin{ejer}\label{matcofej}
    Sea $B = 
    \begin{pmatrix}
        3 & -2 \\
        -1 & 2
    \end{pmatrix}$, calcular $\adj(B)$.   
\end{ejer}

\begin{teo}
Sea $A\in\Mat[\K]{n}{n}$, entonces
\[
    A \adj(A) = \adj(A) A = \det(A) I_{n}.
\]
\end{teo}

\begin{ejem}
    Consideremos la matriz $A$ del ejemplo \ref{matrizcofactores}, se tiene
    \[
        A \adj(A) = \begin{pmatrix}
            0 & 1 & 3 \\
            2 & -1 & 0 \\
            -1 & 2 & -1
        \end{pmatrix}
        \begin{pmatrix}
            1 & 7 & 3 \\
            2 & 3 & 6 \\
            3 & -1 & -2
        \end{pmatrix} = 
        \begin{pmatrix}
            11 & 0 & 0 \\
            0 & 11 & 0 \\
            0 & 0 & 11
        \end{pmatrix} = 11 I_{3}
    \]
    \[
        \adj(A) A = 
        \begin{pmatrix}
            1 & 7 & 3 \\
            2 & 3 & 6 \\
            3 & -1 & -2
        \end{pmatrix}
        \begin{pmatrix}
            0 & 1 & 3 \\
            2 & -1 & 0 \\
            -1 & 2 & -1
        \end{pmatrix} =
        \begin{pmatrix}
            11 & 0 & 0 \\
            0 & 11 & 0 \\
            0 & 0 & 11
        \end{pmatrix} = 11 I_{3}
    \]
    
    \noindent Como el $\det(A) = 11$ (verifique), entonces $A \adj(A) = \adj(A) A = \det(A) I_{3}$.
\end{ejem}

\begin{ejer} 
    Dada la matriz $B$ del ejercicio \ref{matcofej}, compruebe que:
    \[
        B \adj(B) = \adj(B) B = \det(B) I_{2}.
    \]
\end{ejer}

\begin{cor}
Sea $A\in\Mat[\K]{n}{n}$. Si $\det(A) \neq 0$, entonces 
$A$ es invertible y 
\[
    A^{- 1} = \dfrac{1}{\det(A)}\adj(A).
\]
\end{cor}

\begin{ejem}
    Consideremos nuevamente la matriz $A$ del ejemplo \ref{matrizcofactores}.
    
    \noindent Como el $\det(A) \not = 0$, entonces se tiene:
    \[
        A^{-1} = \dfrac{1}{11}
        \begin{pmatrix}
            1 & 7 & 3 \\
            2 & 3 & 6 \\
            3 & -1 & -2
        \end{pmatrix} = 
        \begin{pmatrix}
            \dfrac{1}{11} & \dfrac{7}{11} & \dfrac{3}{11} \\
            \dfrac{2}{11} & \dfrac{3}{11} & \dfrac{6}{11} \\
            \dfrac{3}{11} & -\dfrac{1}{11} & -\dfrac{2}{11}
        \end{pmatrix}
    \]
\end{ejem}

\begin{ejer} Dada la matriz $B$ del ejemplo \ref{matcofej}, calcule su matriz inversa.
\end{ejer}

\begin{teo}\label{TeoInvAdj}
    Sea $A \in \Mat[\K]{n}{n}$. Si $A$ es no singular, entonces $\det(A) \neq 0$ y 
    \[
        \det(A^{-1}) = \dfrac{1}{\det(A)}.
    \]
\end{teo}

\begin{ejem}
    Sea $A$ la matriz del ejemplo \ref{matrizcofactores}. Se tiene que:
    \[
        \det(A^{-1}) = \dfrac{1}{11} = \dfrac{1}{\det(A)}
    \]
\end{ejem}

\begin{ejer} Compruebe la igualdad del teorema \ref{TeoInvAdj} para la matriz del ejemplo \ref{matcofej}.
\end{ejer}

\begin{prop}\label{Pro_NS}
Sea $A\in\Mat[\K]{n}{n}$. Una matriz $A$ es no singular si y sólo si $\det(A)\neq0$.
\end{prop}

\begin{ejer} Hacer la demostración de la proposición \ref{Pro_NS}.
\end{ejer}

\begin{prop}
Sea $A\in\Mat[\K]{n}{n}$. El sistema homogéneo $Ax=0$ tiene solución no trivial si y sólo si $\det(A)=0$.
\end{prop}

\begin{ejer}
Dado el sistema lineal homogéneo:
\[
\systeme[x_1 x_2 x_3]{x_1+5x_2-4x_3=0, x_1-2x_2+x_3=0, 3x_1+x_2-2x_3=0}
\]
Obtenga la solución del sistema.
\end{ejer}
\begin{prop}
Sean $A\in\Mat[\K]{n}{n}$ y $b\in\Mat[\K]{n}{1}$. El sistema $Ax=b$ tiene solución única si y sólo si $\det(A)\neq0$ y su solución es
\[
    A^{- 1}b.
\]
\end{prop}
\begin{ejer}
Dado el sistema lineal:
\[
\systeme[x_1 x_2 x_3]{x_1-4x_2+x_3=5, 2x_1-6x_2-x_3=8, -x_1-5x_3=-12}
\]
Obtenga la solución del sistema.
\end{ejer}
%\begin{teo}[Equivalencias no singulares]
%Sea $A\in\Mat[\K]{n}{n}$, los siguientes enunciados son equivalentes.
%\begin{enumerate}
%\item $A$ es no singular;
%\item el sistema $Ax=0$ tiene solamente la solución trivial;
%\item $A$ es equivalente por filas a $I_{n}$;
%\item $\rang(A)=n$;
%\item el sistema lineal $Ax=b$ tiene una solución única para cada $b\in\Mat[\K]{n}{1}$; y
%\item $\det(A)\neq0$.
%\end{enumerate}
%\end{teo}

%\subsection{Regla de Cramer}
%\begin{teo}[Regla de Cramer]
 %   Sean $A\in\Mat[\K]{n}{n}$ y $b\in\Mat[\K]{n}{1}$. Consideremos el sistema de ecuaciones
  %  \[
   %     \l\{
    %    \begin{array}{c@{\,}c@{\,}c@{\,}c@{\,}c@{\,}c@{\,}c}
     %   a_{11} x_1 & + & a_{12} x_2 & + \, \cdots \, + & a_{1n} x_n & = & b_1, \\
      %  a_{21} x_1 & + & a_{22} x_2 & + \, \cdots \, + & a_{2n} x_n & = & b_2, \\
       % \vdots &    &     \vdots & & \vdots & &  \vdots \\
     %   a_{n1} x_1 & + & a_{n2} x_2 & + \, \cdots \, + & a_{nn} x_n & = & b_n, \\
      %  \end{array}
     %   \r.
    %\]
    %y denotemos por $A$ a su matriz de coeficientes y $b$ a su matriz de constantes. Si $\det(A) \neq 0$, el sistema tiene solución única y está dada por:
  %  \[
   %     x_j = \dfrac{\det(A_j)}{\det(A)},
    %\]
    %para $j\in I$ y donde $A_j$ es la matriz que se obtiene reemplazando la $j$-ésima columna de $A$ por $b$.
%\end{teo}

%\begin{ejem}
 %   Consideremos el sistema de ecuaciones:
  %  \[
   %     \l\{
    %    \begin{array}{c@{\,}c@{\,}c@{\,}c@{\,}c@{\,}c@{\,}c}
   %     2x_1 & + & x_2  & - & x_3  & = & 3 \\
    %    x_1  & + & x_2  & + & x_3  & = & 1 \\
     %   x_1  & - & 2x_2 & - & 3x_3 & = & 4 \\
      %  \end{array}
       % \r.
    %\]
    %En donde:
    %\[
     %   A = 
        %\begin{pmatrix}
            %2 & 1 & -1 \\
            %1 & 1 & 1 \\
            %1 & -2 & -3
%        \end{pmatrix} \qquad y \qquad b =
        %\begin{pmatrix}
         %   3 \\
          %  1 \\
           % 4
        %\end{pmatrix}
    %\]
    %Se sigue que $\det(A) = 5$. Puesto que $\det(A) \not = 0$, el sistema tiene solución única, la cual está determinada por:
  %  \[
   %     x_1 = \dfrac{
        %\begin{vmatrix}
         %   3 & 1 & -1 \\
          %  1 & 1 & 1 \\
           % 4 & -2 & -3
        %\end{vmatrix}}{5} = 2
   % \]
   % \[
    %    x_2 = \dfrac{
        %\begin{vmatrix}
         %   2 & 3 & -1 \\
          %  1 & 1 & 1 \\
           % 1 & 4 & -3
        %\end{vmatrix}}{5} = -1
    %\]
    %\[
     %   x_3 = \dfrac{
        %\begin{vmatrix}
         %   2 & 1 & 3 \\
          %  1 & 1 & 1 \\
           % 1 & -2 & 4
%        \end{vmatrix}}{5} = 0
 %   \]
  %  Así, la solución única del sistema es $x = \begin{pmatrix}2&-1&0\end{pmatrix}^{\intercal}$.
%\end{ejem}

%\begin{ejer}
 %   Resuelva el sistema de ecuaciones:
 %   \[
  %      \l\{
   %     \begin{array}{c@{\,}c@{\,}c@{\,}c@{\,}c}
    %    -ax_1  & + & 2bx_2 & = & c \\
     %   2ax_1 & - & 5bx_2 & = & 3c
      %  \end{array}
     %   \r.
    %\]
    %donde $a$, $b$ y $c$ son números reales, con $a b \not = 0$

  %   \[
   %     \l\{
    %    \begin{array}{c@{\,}c@{\,}c@{\,}c@{\,}c@{\,}c@{\,}c}
     %   3x  & + & 2y & + & z &= & 2 \\
      %  2x  & + & 2y & + & z &= & 2\\
       % x  & + & y & + & z &= & 1
        %\end{array}
        %\r.
    %\]


%\end{ejer}
\end{document}